\documentclass[11pt,a4paper]{article}
\usepackage[utf8]{inputenc}
\usepackage[T1]{fontenc}
\usepackage[portuguese]{babel}
\usepackage{xcolor}
\usepackage{hyperref}
\usepackage{geometry}
\usepackage{tcolorbox}

\geometry{margin=2.5cm}

\definecolor{primary}{RGB}{41, 128, 185}
\definecolor{secondary}{RGB}{44, 62, 80}
\definecolor{codebg}{RGB}{240, 240, 240}

\hypersetup{colorlinks=true, linkcolor=primary, urlcolor=primary}

\title{\color{secondary}\Huge\textbf{Exercício 5.1: Faça você mesmo um CRD e Controller}}
\author{\color{primary}DevOps com Kubernetes -- 2026}
\date{}

\begin{document}
\maketitle

\section*{\color{primary}Instruções do Exercício}
\begin{tcolorbox}[colback=blue!5!white,colframe=blue!75!black,title=Exercício 5.1: Faça você mesmo um CRD e Controller]
Neste exercício, precisamos criar um recurso do zero. Vamos dar o nome de \textit{DummySite} a um recurso que extraia automaticamente o HTML de uma página da web e crie uma cópia do site.

- Crie uma \textit{CustomResourceDefinition} (CRD) definindo a propriedade \texttt{website\_url}.

- Desenvolva um controlador (\textit{controller}) personalizado que obterá e analisará os eventos de todos os \textit{DummySites} criados a partir das APIS REST do Kubernetes.

- Faça o controlador gerar e preencher dinamicamente as requisições do tipo \textit{Job} de que seu código precisará para rodar o scraping e concretizar a clonagem da página.
\end{tcolorbox}

\end{document}
